% !TeX root = ../../Book2.tex
\chapter{Laplace 与 Poisson 方程}
\label{b2:ch:13}
\BookChapterNumber{b2:ch:13}
% 证明义务与实际来源范围见 tmp/book2-part02-20260916/plan.md。
% Axler2001Harmonic 本轮实读pp.4--8、12--18、47--50、144--146。
% 既有第一章Green公式、第五章负Laplace基本解为当前同卷证明输入。
\begin{chapterabstract}
本章研究没有时间变量的平衡方程。调和函数的均值结构产生最大值原理、Harnack 比较与内部光滑性；Green 恒等式则把内部源量与边界通量联系起来。利用既有基本解构造 Newton 势后，进一步分析奇核的一阶可积性和二阶抵消，说明右端正则性怎样影响经典解。最后通过球面反演与平面反射构造 Green 函数，完成球和半空间上的连续 Dirichlet 问题，并明确一般区域与无穷远条件所留下的边界。
\end{chapterabstract}

\begin{learninggoals}
\begin{itemize}
\item\label{b2:item:13-goal-mean}从球面和球内平均理解调和性，掌握最大值、Harnack 链和 Liouville 结论的假设。
\item\label{b2:item:13-goal-boundary}使用外法向 Green 恒等式检查 Dirichlet、Neumann 和 Robin 数据的唯一性及相容条件。
\item\label{b2:item:13-goal-potential}构造 Newton 势，区别有界、连续与 Hölder 密度所能保证的正则性。
\item\label{b2:item:13-goal-green}核对 Green 函数的点源、零边界、正性和对称性，区分表示公式与存在性证明。
\item\label{b2:item:13-goal-models}用显式 Poisson 核解球与半空间的问题，说明镜像系数、维数和增长类各起什么作用。
\end{itemize}
\end{learninggoals}

本章沿用\chapref{b2:ch:01}的向量分析与外法向约定，以及\chapref{b2:ch:05}已经证明的负 Laplace 基本解。方程统一写为 \(-\Delta u=f\)，所以正源与零 Dirichlet 数据对应非负解。\chapref{b2:ch:10}的 Green 核首先服务于有限区间；这里把相同的逆算子思想放入欧氏空间的几何边界中。

\ProseParagraphBreak

\secref{b2:sec:1301}的局部调和理论先独立建立，Harnack 的证明不预借后面才出现的 Poisson 核。\secref{b2:sec:1302}整理边界与能量接口，\secref{b2:sec:1303}处理体积源。\secref{b2:sec:1304}在明确假设下一般讨论 Green 表示，并在球与半空间完成存在性；任意粗糙域上的边界正则性不在这组显式构造中自动得到。

\ProseParagraphBreak

Sheldon Axler、Paul Bourdon 与 Wade Ramey 的《Harmonic Function Theory》\cite[Chapters 1, 3, and 7]{Axler2001Harmonic}可伴读均值、正调和函数及半空间核，比较其单位球面概率测度与本书未归一化面积测度的差别。Lawrence C. Evans 的《Partial Differential Equations》\cite{Evans2010Partial}可用于继续学习一般椭圆方程与弱解；本章不以未展开的一般正则性结论替代显式核的计算。

% !TeX root = ../../Book2.tex
\section{调和函数}
\label{b2:sec:1301}

热传导达到平衡而区域内部没有热源时，温度应当满足 Laplace 方程。这样的函数在一个小球中心的值由球面或球内平均决定，局部不能凭空出现新的峰值。这种平均结构还会产生光滑性、比较估计和唯一性，远比单独观察一个二阶微分等式更强。

\Needspace{6\baselineskip}
\begin{definition}\label{b2:def:13-harmonic}
设整数 \(d\geq1\)，\(\Omega\subset\mathbb R^d\) 为开集，\(\mathbb K\in\{\mathbb R,\mathbb C\}\)。若 \(u\in C^2(\Omega;\mathbb K)\) 且 \(\Delta u=0\)，则称 \(u\) 为 \(\Omega\) 上的\term[diaohe-hanshu]{调和函数}[harmonic function]。除另作说明外，本章先取实值函数；复值函数的调和性按实部和虚部分别理解。
\end{definition}

一次函数调和，二次函数 \(x_1^2-x_2^2\) 也调和，但 \(|x|^2\) 的 Laplace 等于 \(2d\)，并不调和。一维调和函数恰为仿射函数；本章涉及球面核和基本解的主要计算固定 \(d\geq2\)。沿用 \(\omega_d=|B(0,1)|\)、\(\sigma_{d-1}=|\partial B(0,1)|=d\omega_d\)，分别表示体积与表面积。

\subsection{平均值性质}
\label{b2:subsec:130101}

\Needspace{6\baselineskip}
\begin{theorem}[调和函数的平均值性质]\label{b2:thm:13-mean-value}
设整数 \(d\geq1\)，\(\Omega\subset\mathbb R^d\) 开、\(a\in\Omega\)、\(R>0\)。若实值或复值函数 \(u\) 在 \(\Omega\) 上调和，且 \(\overline{B(a,R)}\subset\Omega\)，则对 \(0<r\leq R\)，有
\[
 u(a)=\dfrac1{\sigma_{d-1}r^{d-1}}\int_{\partial B(a,r)}u\,\mathrm{d}S
 =\dfrac1{\omega_dr^d}\int_{B(a,r)}u\,\mathrm{d}x.
\]
\end{theorem}
\begin{proof}
将球面平均写到固定单位球面上，记
\[
 M(r)=\dfrac1{\sigma_{d-1}}\int_{\partial B(0,1)}u(a+r\theta)\,\mathrm{d}S_\theta.
\]
对固定的紧半径区间可在积分下求导。换回半径为 \(r\) 的球面，再用散度定理，得到
\[
 M'(r)=\dfrac1{\sigma_{d-1}r^{d-1}}
          \int_{\partial B(a,r)}\partial_\nu u\,\mathrm{d}S
       =\dfrac1{\sigma_{d-1}r^{d-1}}\int_{B(a,r)}\Delta u\,\mathrm{d}x=0.
\]
所以 \(M\) 为常数，而 \(r\downarrow0\) 时由连续性有 \(M(r)\to u(a)\)，得到球面公式。最后用极坐标沿 \(0<s<r\) 积分，球内积分等于 \(\int_0^r\sigma_{d-1}s^{d-1}u(a)\,\mathrm{d}s=\omega_dr^du(a)\)。
\end{proof}

这个\term[pingjunzhi-xingzhi]{平均值性质}[mean-value property]不含导数，因而也可以反过来用于识别调和性。

\Needspace{6\baselineskip}
\begin{proposition}[平均值性质的逆命题]\label{b2:prop:13-converse-mean}
设整数 \(d\geq1\)，\(\Omega\subset\mathbb R^d\) 为开集，\(\mathbb K\in\{\mathbb R,\mathbb C\}\)，\(u\in C(\Omega;\mathbb K)\)。若每个闭球包含于 \(\Omega\) 时，球面平均均等于中心值，则 \(u\) 光滑且调和。若 \(u\) 连续并在分布意义下满足 \(\Delta u=0\)，也有同样结论。
\end{proposition}
\begin{proof}
取非负径向磨光函数 \(\rho\in C_c^\infty(B(0,1))\)，积分为 \(1\)，并令 \(\rho_\varepsilon(x)=\varepsilon^{-d}\rho(x/\varepsilon)\)。球面平均公式按径向权积分，给出 \(u*\rho_\varepsilon=u\)，只要核的支集不触及边界。对任一闭包包含于 \(\Omega\) 的球，可取同一个小 \(\varepsilon\) 使这一等式在该球附近成立。左侧光滑，故 \(u\) 局部光滑。

现在对 \(u(a+r\theta)\) 作二阶 Taylor 展开。球面上 \(\theta_j\) 的平均为零，\(\theta_i\theta_j\) 在 \(i\ne j\) 时平均为零，\(\theta_j^2\) 的平均为 \(1/d\)。因此球面平均为
\[
 u(a)+\dfrac{r^2}{2d}\Delta u(a)+o(r^2).
\]
它恒等于 \(u(a)\)，除以 \(r^2\) 再令 \(r\downarrow0\) 可得 \(\Delta u(a)=0\)。

若起初只知分布方程，内部磨光 \(u_\varepsilon=u*\rho_\varepsilon\) 满足经典方程 \(\Delta u_\varepsilon=0\)，因为分布求导与磨光交换。连续性保证 \(u_\varepsilon\to u\) 在紧集上一致成立。对任意固定内部闭球，令 \(\varepsilon\downarrow0\)，便将 \(u_\varepsilon\) 的球面平均公式传给 \(u\)，再应用前半部分。
\end{proof}

球面平均与球内平均的归一化不可混用。若将球面的未归一化面积积分直接视为平均，会额外带入 \(\sigma_{d-1}r^{d-1}\)。同样，平均值性质针对整个小球而非某个坐标方向；只要求每条坐标线上的中点值等于端点平均，得到的是另一种更强且不同的条件。

有时只知道每个中心附近的平均性质，允许使用的半径还随中心变化。这比刚才对全部内部球的假设弱，不能直接为整个邻域选定同一个卷积尺度。

\Needspace{6\baselineskip}
\begin{proposition}[局部平均值性质的逆命题]\label{b2:prop:13-local-converse-mean}
设整数 \(d\geq1\)，\(\Omega\subset\mathbb R^d\) 为开集，\(\mathbb K\in\{\mathbb R,\mathbb C\}\)，\(u\in C(\Omega;\mathbb K)\)。若对每个 \(a\in\Omega\)，存在 \(r_a>0\)，使 \(\overline{B(a,r_a)}\subset\Omega\)，并且对所有 \(0<r<r_a\) 有
\[
 u(a)=\frac1{\sigma_{d-1}r^{d-1}}
           \int_{\partial B(a,r)}u\,\mathrm{d}S,
 \qquad \sigma_{d-1}=|S^{d-1}|,
\]
则 \(u\) 在 \(\Omega\) 内光滑且调和。
\end{proposition}

证明延后至\subsecref{b2:subsec:13-ball-boundary}，在那里先建立球上连续 Dirichlet 问题的解，再用它将局部均值性质推广到全部内部球。本节此前已经证明的\cref{b2:prop:13-converse-mean}不需要这一步构造。

\subsection{最大值原理}
\label{b2:subsec:130102}

\Needspace{6\baselineskip}
\begin{theorem}[调和函数的强最大值原理]\label{b2:thm:13-strong-maximum}
设整数 \(d\geq1\)，\(\Omega\subset\mathbb R^d\) 为连通开集，\(u\) 为 \(\Omega\) 上的实值调和函数。若 \(u\) 在 \(\Omega\) 内达到它在该区域上的最大值或最小值，则 \(u\) 为常数。
\end{theorem}
\begin{proof}
设最大值 \(M\) 在 \(a\) 处达到。对一个内部小球，非负连续函数 \(M-u\) 的球内平均为零，所以它在整个球上为零：若某点取正值，连续性会给出一个正测度区域，其积分为正。故集合 \(\{u=M\}\) 在 \(\Omega\) 中开；由连续性它也闭，且非空。连通性使它等于 \(\Omega\)。最小值情形对 \(-u\) 使用同一论证。
\end{proof}

这项\term[diaohe-qiang-zuidazhi-yuanli]{调和函数强最大值原理}[strong maximum principle for harmonic functions]中的“达到”不能改成只有有限上确界。例如 \(u(x_1,x_2)=e^{x_1}\cos x_2\) 在半平面 \(x_1<0\) 上调和且有界，上确界为 \(1\)，却不在内部达到。连通性也不可省去，否则不同连通分支可以取不同常数。

\Needspace{6\baselineskip}
\begin{corollary}[边界最大值与比较]\label{b2:cor:13-boundary-maximum}
设整数 \(d\geq1\)，\(\Omega\subset\mathbb R^d\) 为非空有界开集，实值 \(u\in C(\overline\Omega)\) 在 \(\Omega\) 内调和，则
\[
 \min_{\partial\Omega}u\leq u(x)\leq\max_{\partial\Omega}u.
\]
若两个这样的调和函数在边界满足 \(u\leq v\)，则区域内也满足 \(u\leq v\)。
\end{corollary}
\begin{proof}
连续函数在紧集 \(\overline\Omega\) 上达到最大值。若它严格超过边界最大值，只能在某个内部连通分支达到；强最大值原理使函数在该分支恒定，沿分支内点趋于其边界便与边界条件矛盾。对 \(-u\) 得到下界，对 \(u-v\) 得到比较结论。这里不需要边界光滑。
\end{proof}

复值调和函数也有模的最大值原理。若 \(|u|\leq M\) 且 \(|u(a)|=M\)，选单位复数 \(c\) 使 \(cu(a)=M\)。实值调和函数 \(\operatorname{Re}(cu)\leq M\) 在内部达到最大值，故恒为 \(M\)；再由 \(|cu|\leq M\) 得虚部为零。于是 \(u\) 为常数。实函数的大小比较和复函数的模比较应分别使用，不能把复数直接放进大小不等式。

\subsection{Harnack 不等式}
\label{b2:subsec:130103}

最大值原理比较内部和边界。若函数非负，则平均值性质还可以比较两个内部点，而无需事先知道边界数据。

\Needspace{6\baselineskip}
\begin{lemma}\label{b2:lem:13-local-harnack}
设整数 \(d\geq1\)，\(a\in\mathbb R^d\)、\(R>0\)，实值 \(u\ge0\) 在 \(B(a,R)\) 上调和。当 \(|x-a|\leq R/4\) 时，
\[
 2^{-d}u(a)\leq u(x)\leq(3/2)^d u(a).
\]
\end{lemma}
\begin{proof}
包含关系 \(B(a,R/4)\subset B(x,R/2)\subset B(a,3R/4)\) 与非负性给出两个积分比较。对三只球分别使用平均值性质，再除以 \(|B(x,R/2)|\)，即得体积比 \(2^{-d}\) 与 \((3/2)^d\)。所有闭球都留在调和域内部，端点情形可由稍缩半径后取极限得到。
\end{proof}

\Needspace{6\baselineskip}
\begin{theorem}[Harnack]\label{b2:thm:13-harnack}
设整数 \(d\geq1\)，\(a\in\mathbb R^d\)、\(R>0\)。若实值 \(u\ge0\) 在 \(B(a,R)\) 上调和，\(0<r<R\)，则存在只依赖 \(d,r/R\) 的常数 \(C\)，使
\[
 \sup_{B(a,r)}u\leq C\inf_{B(a,r)}u.
\]
更一般地，若 \(\Omega\subset\mathbb R^d\) 为连通开集、实值 \(u\ge0\) 在 \(\Omega\) 上调和、\(K\Subset\Omega\)，则存在只依赖 \(K,\Omega\) 的相同类型比较常数。
\end{theorem}
\begin{proof}
连接内球中两点的线段完全留在内球。令 \(\delta=R-r\)，将线段分成长度至多 \(\delta/4\) 的有限段；每个分点为中心的半径 \(\delta\) 开球都在原域内。逐段使用局部引理，一步的双向比较均可用 \(2^d\) 控制，段数至多 \(\lceil8r/\delta\rceil+1\)。连乘即得只依赖半径比的常数，最后对两点分别取上、下确界。

一般连通开集中的两点可以用域内折线连接：从一个固定点可用折线到达的集合在域中既开又闭，因此是全域。每条有限折线与边界保持正距离，可以被有限条上述小球链覆盖。对紧集 \(K\)，再选有限个终端小球覆盖它，合并有限条连接链，便得到对全部 \(K\) 统一的常数。比较只用于非负函数，故每次积分包含关系都保留方向。
\end{proof}

这称为\term[harnack-budengshi]{Harnack 不等式}[Harnack's inequality]。非负条件是实质性的：函数 \(u(x)=x_1\) 调和，但一个包含原点的球中同时有正、负值，不满足这种正比例比较。对非负调和函数，如果在一点取零，Harnack 或强最小值原理都会使它在整个连通域为零。

\Needspace{6\baselineskip}
\begin{corollary}[Harnack 收敛原理]\label{b2:cor:13-harnack-convergence}
设整数 \(d\geq1\)，\(\Omega\subset\mathbb R^d\) 为连通开集，\((u_n)\) 为 \(\Omega\) 上逐点递增的实值调和函数列，\(a\in\Omega\)。若 \(u_n(a)\) 在某一点有有限上界，则 \(u_n\) 在每个紧集上一致收敛到调和函数；若该点的值趋于 \(+\infty\)，则在每个紧集上一致趋于 \(+\infty\)。
\end{corollary}
\begin{proof}
第一种情形中，对 \(m>n\) 将 Harnack 应用于非负函数 \(u_m-u_n\)。它在紧集上的上确界由常数乘 \(u_m(a)-u_n(a)\) 控制，故为一致 Cauchy 列。极限连续，球面平均公式按局部一致极限保留，由已证明的\cref{b2:prop:13-converse-mean}知极限调和。第二种情形对 \(u_n-u_1\) 作反向的紧集比较，其下确界随 \(u_n(a)-u_1(a)\) 趋于无穷。
\end{proof}

\subsection{内部正则性与 Liouville 定理}
\label{b2:subsec:130104}

\Needspace{6\baselineskip}
\begin{proposition}[内部导数估计]\label{b2:prop:13-interior-derivatives}
设整数 \(d\geq1\)，\(a\in\mathbb R^d\)、\(R>0\)，实值或复值 \(u\) 在 \(B(a,2R)\) 上调和，则对每个多重指标 \(\alpha\)，
\[
 \sup_{B(a,R)}|\partial^\alpha u|
 \leq C_{d,\alpha}R^{-d-|\alpha|}\int_{B(a,2R)}|u|\,\mathrm{d}x.
\]
\end{proposition}
\begin{proof}
先将球面平均按径向磨光权积分。对 \(x\in B(a,R)\)，有 \(u(x)=\int\rho_R(x-y)u(y)\,\mathrm{d}y\)，核支集限于 \(B(a,2R)\)。将导数移到光滑核，\(\|\partial^\alpha\rho_R\|_\infty=R^{-d-|\alpha|}\|\partial^\alpha\rho\|_\infty\)，得到所求估计。等式本身也再次证明了调和函数的任意阶光滑性。
\end{proof}

\Needspace{6\baselineskip}
\begin{proposition}[局部可积调和函数的代表]\label{b2:prop:13-harmonic-l1-representative}
设整数 \(d\geq1\)，\(\Omega\subset\mathbb R^d\) 为开集，\(\mathbb K\in\{\mathbb R,\mathbb C\}\)。若 \(h\in L^1_{\mathrm{loc}}(\Omega;\mathbb K)\)，且 \(\Delta h=0\) 在分布意义下成立，则 \(h\) 有唯一的光滑调和代表。因而内部导数估计也适用于这个代表，右端仍使用原等价类的绝对积分。
\end{proposition}
\begin{proof}
固定 \(\overline{B(a,3R)}\subset\Omega\)，用半径 \(\varepsilon<R/2\) 的磨光核构造 \(h_\varepsilon\)。分布求导与卷积交换，所以 \(h_\varepsilon\) 在 \(B(a,2R)\) 上光滑且调和。局部的 \(L^1\) 逼近恒等给出 \(h_\varepsilon\to h\) 于 \(L^1(B(a,2R))\)。对两个磨光尺度之差使用内部估计，得到
\[
 \sup_{B(a,R)}|\partial^\alpha(h_\varepsilon-h_\delta)|
 \leq C_{d,\alpha}R^{-d-|\alpha|}
       \|h_\varepsilon-h_\delta\|_{L^1(B(a,2R))}\longrightarrow0.
\]
因此任意有限阶导数都在内部小球上一致 Cauchy；逐次使用微积分基本定理，极限是光滑函数，其各阶导数正是相应的一致极限。二阶导数收敛使极限仍调和，零阶 \(L^1\) 收敛又将它识别为 \(h\) 的几乎处处代表。

这些内部球覆盖 \(\Omega\)。两个局部代表在重叠处几乎处处相同，而连续函数若在一个开集上几乎处处为零，便处处为零。因此局部代表能够拼合，且光滑代表唯一。最后对任一内部球使用前一命题，代表与 \(h\) 的绝对积分相同，故估计保留。
\end{proof}

这一步把分布方程与经典调和性连接起来，所用条件只是局部可积性。磨光函数的每个导数虽各自存在，仅有 \(L^1\) 收敛却原本不足以传递导数；真正提供一致控制的是调和函数的内部估计。

\Needspace{6\baselineskip}
\begin{theorem}[Liouville]\label{b2:thm:13-liouville}
设整数 \(d\geq1\)，\(\mathbb K\in\{\mathbb R,\mathbb C\}\)，\(u\) 为整个 \(\mathbb R^d\) 上的 \(\mathbb K\) 值调和函数。若 \(u\) 全局有界，则 \(u\) 为常数。实值 \(u\) 只要有全局上界或下界，也必为常数。若存在 \(C>0\) 和非负整数 \(m\)，使全部 \(x\in\mathbb R^d\) 满足 \(|u(x)|\leq C(1+|x|)^m\)，则 \(u\) 为次数不超过 \(m\) 的调和多项式。
\end{theorem}
\begin{proof}
有界情形对任意点使用一阶内部估计，让 \(R\to\infty\)，得到梯度处处为零。若实函数上有界，取 \(v=M-u\geq0\)。内部估计中的绝对积分等于 \(v\) 的积分，而平均值性质给出 \(\int_{B(a,2R)}v=\omega_d(2R)^dv(a)\)。因此 \(|\nabla v(a)|\leq C_dR^{-1}v(a)\to0\)。下有界时改用 \(u-M\)。

在多项式增长条件下，固定 \(a\)，取 \(|\alpha|=m+1\)。内部估计的右侧为 \(O(R^{-1})\)，故全部 \(m+1\) 阶偏导数为零。对连接原点与任意点的线段使用一元 Taylor 公式，余项中的 \(m+1\) 阶方向导数为零，所以 \(u\) 恰等于在原点的 \(m\) 阶 Taylor 多项式。该多项式仍满足原来的调和方程。
\end{proof}

上述结论统称调和函数的\term[diaohe-liouville]{Liouville 定理}[Liouville's theorem for harmonic functions]。增长条件不能删除，例如 \(x_1\) 是非常数整调和函数；仅有界于某个真子域时也不能套用全空间结论。后面处理半空间唯一性时，必须同时利用边界条件和适当的增长类。

% !TeX root = ../../Book2.tex
\section{Green 恒等式}
\label{b2:sec:1302}

均值性质适合局部比较，Green 恒等式则把区域内部的微分与边界通量放在同一个式子中。我们沿用\cref{b2:cor:01-vector-green}，把外法向固定为 \(\nu\)，并将 Poisson 方程统一写成 \(-\Delta u=f\)。本节能量与边界积分先对有界 \(C^1\) 区域，以及在其闭包某邻域内为 \(C^2\) 的函数推导；只用最大值原理的结论会另行放宽边界条件。

\subsection{第一 Green 恒等式}
\label{b2:subsec:130201}

在当前符号下，第一 Green 恒等式为
\begin{equation}\label{b2:eq:13-green-first}
 \int_\Omega v\cdot(-\Delta u)\,\mathrm{d}x
 =\int_\Omega\nabla u\cdot\nabla v\,\mathrm{d}x
  -\int_{\partial\Omega}v\cdot\partial_\nu u\,\mathrm{d}S.
\end{equation}
它是\eqref{b2:eq:01-vector-green-first}的移项形式。若 \(-\Delta u=f\)，取 \(v=u\) 得到
\[
 \int_\Omega|\nabla u|^2\,\mathrm{d}x
 =\int_\Omega f\cdot u\,\mathrm{d}x
  +\int_{\partial\Omega}u\cdot\partial_\nu u\,\mathrm{d}S.
\]
边界上给定函数值与给定法向导数，控制的是同一边界乘积的不同因子。它们不应被当作可以任意同时指定的两套数据。

\ProseParagraphBreak

例如在球 \(B(0,R)\) 上，函数 \(u(x)=(R^2-|x|^2)/(2d)\) 满足 \(-\Delta u=1\)、\(u|_{\partial B}=0\)。能量等式此时化为 \(\int|\nabla u|^2=\int u\)。左侧由 \(\nabla u=-x/d\) 计算，右侧由极坐标计算，二者都等于 \(\omega_dR^{d+2}/(d(d+2))\)。边界值为零消去了通量乘积，并不意味着法向导数本身为零；事实上 \(\partial_\nu u=-R/d\)。

\begin{proposition}[Dirichlet 能量极小性]\label{b2:prop:13-dirichlet-energy}
设 \(u\) 为实值调和函数，\(v\) 与 \(u\) 具有相同边界值，且两者满足本节的光滑假设。则
\[
 \int_\Omega|\nabla v|^2\,\mathrm{d}x
 =\int_\Omega|\nabla u|^2\,\mathrm{d}x
  +\int_\Omega|\nabla(v-u)|^2\,\mathrm{d}x.
\]
因此 \(u\) 在同边界值的光滑函数中使梯度能量最小。
\end{proposition}
\begin{proof}
令 \(w=v-u\)。展开 \(|\nabla(u+w)|^2\) 后，交叉项为 \(2\int\nabla u\cdot\nabla w\)。由第一 Green 恒等式，\(\Delta u=0\) 及 \(w=0\) 于边界，该项为零，得到恒等式。
\end{proof}

这只证明已有调和解的极小性，没有仅凭能量非负就证明极小元存在。要从极小化序列构造解，还需选择完备的函数空间并处理紧性与边界约束；这些问题在本卷的变分法与弱解部分继续研究。

\subsection{第二 Green 恒等式}
\label{b2:subsec:130202}

对 \(u,v\) 交换第一恒等式中的位置再相减，得到
\begin{equation}\label{b2:eq:13-green-second}
 \int_\Omega\bigl(u\cdot\Delta v-v\cdot\Delta u\bigr)\,\mathrm{d}x
 =\int_{\partial\Omega}
       \bigl(u\cdot\partial_\nu v-v\cdot\partial_\nu u\bigr)\,\mathrm{d}S.
\end{equation}
如果两者都调和，边界上的两种通量配对相等。这个对称性是后面 Green 函数对称性的来源。

\begin{proposition}[Neumann 数据的相容条件]\label{b2:prop:13-neumann-compatibility}
若 \(-\Delta u=f\) 于 \(\Omega\)，且 \(\partial_\nu u=g\) 于 \(\partial\Omega\)，则必有
\[
 \int_\Omega f\,\mathrm{d}x+\int_{\partial\Omega}g\,\mathrm{d}S=0.
\]
\end{proposition}
\begin{proof}
对 \(\Delta u\) 使用散度定理，\(\int_\Omega\Delta u=\int_{\partial\Omega}\partial_\nu u\)。代入方程和边界值，再移项即得。
\end{proof}

对上一小节的球内例子，\(f=1\)、\(g=-R/d\)，两个积分为 \(\omega_dR^d\) 和 \(-\sigma_{d-1}R^d/d\)，正好抵消。若同样取 \(f=1\) 却要求整个边界上 \(g=0\)，便不可能有解。相容条件是对总源量与总通量的约束；即使它成立，本节也不据此直接宣布任意边界上的存在性。

\ProseParagraphBreak

在复值情形，\eqref{b2:eq:13-green-second}仍可作为双线性恒等式使用；讨论自伴或正能量时则应把一个函数换成共轭。比如 \(\int\nabla u\cdot\nabla\overline u=\int|\nabla u|^2\)，而不加共轭的平方和可能不是非负实数。

\subsection{唯一性}
\label{b2:subsec:130203}

\begin{theorem}[三类边界条件下的唯一性]\label{b2:thm:13-boundary-uniqueness}
固定有界连通 \(C^1\) 区域和给定的 \(f\)。在本节的光滑类中，方程 \(-\Delta u=f\) 满足下列结论。
\begin{enumerate}
\item\label{b2:item:13-unique-dirichlet}给定 \(u|_{\partial\Omega}\) 时，解至多一个。
\item\label{b2:item:13-unique-neumann}给定 \(\partial_\nu u\) 时，解在相差常数的意义下至多一个；再给定 \(\int_\Omega u\) 可得到唯一性。
\item\label{b2:item:13-unique-robin}若 \(\beta\) 连续、\(\beta\geq0\) 且 \(\int_{\partial\Omega}\beta\,\mathrm{d}S>0\)，给定 \(\partial_\nu u+\beta u\) 时，解至多一个。
\end{enumerate}
\end{theorem}
\begin{proof}
取两个解之差 \(w\)，则 \(\Delta w=0\)。在 Dirichlet 或 Neumann 情形，第一 Green 恒等式给出 \(\int|\nabla w|^2=0\)，因为边界乘积分别由 \(w=0\) 或 \(\partial_\nu w=0\) 消去。于是 \(\nabla w=0\)，连通性使 \(w\) 为常数；Dirichlet 数据使该常数为零，Neumann 数据则允许常数自由度。

Robin 情形中，\(\partial_\nu w=-\beta w\)，故
\[
 \int_\Omega|\nabla w|^2\,\mathrm{d}x
 +\int_{\partial\Omega}\beta|w|^2\,\mathrm{d}S=0.
\]
两项非负，第一项使 \(w=c\)，第二项变为 \(|c|^2\int\beta=0\)，由严格正积分得 \(c=0\)。复值解取共轭后同样证明。
\end{proof}

Dirichlet 唯一性还可在任意有界开集上由最大值原理取得，只需解属于 \(C^2(\Omega)\cap C(\overline\Omega)\)，无需边界光滑。Neumann 与 Robin 的能量论证则实际使用了边界法向和 Green 恒等式，不能不加说明地移到任意粗糙边界。

\ProseParagraphBreak

Robin 系数的符号也不是装饰。在球上取 \(w(x)=x_1\)，则 \(\Delta w=0\)、\(\partial_\nu w=w/R\)。因此 \(\beta=-1/R\) 时，非零函数 \(w\) 满足齐次 Robin 数据，唯一性失败。若区域不连通，Neumann 情形则允许每个连通分支各加一个常数，只固定整个区域的平均值通常不能消除全部自由度。

\begin{proposition}\label{b2:prop:13-subharmonic-comparison}
设 \(\Omega\) 为有界开集，\(w\in C^2(\Omega)\cap C(\overline\Omega)\)，且 \(\Delta w\geq0\)。则 \(w(x)\leq\max_{\partial\Omega}w\)。因而若 \(-\Delta u\leq-\Delta v\) 且 \(u\leq v\) 于边界，则内部也有 \(u\leq v\)。
\end{proposition}
\begin{proof}
令 \(w_\varepsilon=w+\varepsilon|x|^2\)，则 \(\Delta w_\varepsilon\geq2d\varepsilon>0\)。如果其最大值在内部达到，Hesse 矩阵在该点半负定，会使 Laplace 不大于零，矛盾。因此
\[
 w(x)+\varepsilon|x|^2
 \leq\max_{\partial\Omega}w+\varepsilon\max_{\partial\Omega}|x|^2.
\]
让 \(\varepsilon\downarrow0\) 即得第一项；第二项用于 \(w=u-v\)。
\end{proof}

\begin{corollary}[边界数据的一致稳定性]\label{b2:cor:13-dirichlet-stability}
设 \(u,v\in C^2(\Omega)\cap C(\overline\Omega)\) 在有界区域中满足相同的 Poisson 右端，则
\[
 \|u-v\|_{C(\overline\Omega)}\leq\|u-v\|_{C(\partial\Omega)}.
\]
\end{corollary}
\begin{proof}
差 \(u-v\) 调和，分别对它和负号使用边界最大值原理；复值时使用模的最大值原理。
\end{proof}

这里稳定性只比较相同的内部源项。如果右端也改变，需要另估计由该差产生的势；不能把一个有源方程的解直接当作调和函数使用。

\subsection{边界积分}
\label{b2:subsec:130204}

现在将第二 Green 恒等式中的一个函数换成带极点的基本解。沿用\cref{b2:thm:05-laplace-fundamental}的 \(\Phi_d\)，满足 \(-\Delta\Phi_d=\delta_0\)。极点附近的边界项会保留下一个点值，不能把基本解当作整个区域内的普通调和函数。

\begin{theorem}[Green 表示公式]\label{b2:thm:13-green-representation}
设 \(d\geq2\)，\(u\) 在有界 \(C^1\) 区域 \(\overline\Omega\) 的某邻域内为 \(C^2\)，并记 \(f=-\Delta u\)。则对每个 \(x\in\Omega\)，
\begin{equation}\label{b2:eq:13-green-representation}
 \begin{aligned}
 u(x)={}&\int_\Omega\Phi_d(x-y)\cdot f(y)\,\mathrm{d}y\\
 &+\int_{\partial\Omega}
 \left(\Phi_d(x-y)\cdot\partial_\nu u(y)
       -u(y)\cdot\partial_{\nu_y}\Phi_d(x-y)\right)\,\mathrm{d}S_y.
 \end{aligned}
\end{equation}
\end{theorem}
\begin{proof}
选 \(\varepsilon>0\) 使闭球 \(\overline{B(x,\varepsilon)}\subset\Omega\)，在挖掉这只球的区域上对 \(u(y)\) 和 \(\Phi_d(x-y)\) 使用第二 Green 恒等式。该区域内基本解调和，故
\[
 -\int_{\Omega\setminus\overline{B(x,\varepsilon)}}\Phi_d(x-y)f(y)\,\mathrm{d}y
 =I_{\partial\Omega}+I_\varepsilon,
\]
其中 \(I_{\partial\Omega}\) 是所求公式的边界积分，内球边界项为
\[
 I_\varepsilon=\int_{|y-x|=\varepsilon}
 \left(\Phi_d(x-y)\cdot\partial_\nu u(y)
 -u(y)\cdot\partial_\nu\Phi_d(x-y)\right)\,\mathrm{d}S_y.
\]
内边界相对于挖孔区域的外法向朝向 \(x\)，所以 \(\partial_\nu\Phi_d=1/(\sigma_{d-1}\varepsilon^{d-1})\)。第二项为 \(u\) 的负球面平均，趋于 \(-u(x)\)。第一项的模至多为常数乘 \(\varepsilon\)（\(d\geq3\)）或 \(\varepsilon|\log\varepsilon|\)（\(d=2\)），故趋零。体积分由 \(f\) 有界和 \(\Phi_d\) 局部可积而收敛，令 \(\varepsilon\downarrow0\) 并移项即得。
\end{proof}

表示公式同时出现边界值与法向导数，所以还不是只用 Dirichlet 数据的求解公式。Green 函数的作用正是修改基本解，使它在边界为零，从而消去不希望出现的法向导数项。下一节先处理体积项及其正则性，随后再作这一边界修正。

% !TeX root = ../../Book2.tex
\section{基本解与 Newton 势}
\label{b2:sec:1303}

基本解描述集中于一点的单位源。将它与分布在整个区域中的密度叠加，就得到 Poisson 方程的一个特解。奇核的可积性决定叠加是否存在，核的导数及其抵消则决定这个特解有多光滑。

\subsection{Laplace 算子的基本解}
\label{b2:subsec:130301}

\Cref{b2:thm:05-laplace-fundamental}已经证明，负 Laplace 算子的基本解可取
\[
 \Phi_d(z)=\dfrac{|z|^{2-d}}{(d-2)\sigma_{d-1}}\quad(d\geq3),
 \quad \Phi_2(z)=-\dfrac1{2\pi}\log|z|,
\]
一维则为 \(\Phi_1(z)=-|z|/2\)。它们在分布意义下满足 \(-\Delta\Phi_d=\delta_0\)，在原点外调和。以下涉及奇核积分时固定 \(d\geq2\)。

\ProseParagraphBreak

原点附近的奇性仍局部可积：高维核与体积元相乘后只剩 \(r\,\mathrm{d}r\)，二维对数也可积。但二阶导数的尺度为 \(|z|^{-d}\)，绝对积分已处于临界发散情形。因此“基本解可积”并不能自动许可两次在积分下求导，下面会保留所需的抵消项。

\subsection{Poisson 方程}
\label{b2:subsec:130302}

对 \(f\in C_c^\infty(\mathbb R^d)\)，\cref{b2:prop:05-fundamental-convolution}给出光滑特解 \(u=\Phi_d*f\)。这个构造把导数移到右端的紧支集光滑函数上，得到 \(-\Delta u=f\)。它先解全空间方程；若目标区域上另给边界值，还需加一个调和修正。

\begin{proposition}\label{b2:prop:13-poisson-affine-space}
若 \(u_0\in C^2(\Omega)\) 满足 \(-\Delta u_0=f\)，则全部 \(C^2\) 解恰为 \(u=u_0+h\)，其中 \(h\) 在 \(\Omega\) 上调和。
\end{proposition}
\begin{proof}
线性性使任意两个解之差满足 \(\Delta h=0\)；反过来，给特解加调和函数不改变方程右端。
\end{proof}

这说明基本解不负责选择边界条件。例如球上 \(-\Delta u=1\) 的一个多项式特解为 \(-|x|^2/(2d)\)；若要求边界值为零，再加常数 \(R^2/(2d)\)。非球形区域上需要的修正一般不再是常数，不能沿用球上公式只替换一个“半径”。

\subsection{Newton 势}
\label{b2:subsec:130303}

\begin{definition}\label{b2:def:13-newton-potential}
设 \(f\) 是有界紧支集的可测函数。定义
\[
 (Nf)(x)\coloneqq\int_{\mathbb R^d}\Phi_d(x-y)\cdot f(y)\,\mathrm{d}y,
\]
称为密度 \(f\) 的\term[newton-shi]{Newton 势}[Newtonian potential]。二维沿用同一记号，但其核为对数核。
\end{definition}
\symidx{\(Nf\)：Newton 势}

\begin{proposition}\label{b2:prop:13-newton-c1}
上述积分在每个 \(x\) 绝对收敛，\(Nf\in C^1(\mathbb R^d)\)，并且
\[
 \partial_jNf(x)=\int\partial_j\Phi_d(x-y)\cdot f(y)\,\mathrm{d}y,
 \quad-\Delta Nf=f
\]
后一等式取分布意义。
\end{proposition}
\begin{proof}
固定 \(x\)。在 \(y=x\) 附近用 \(f\) 的本性界和 \(\Phi_d\) 的局部可积性，远处则用 \(f\) 的紧支集，故积分绝对收敛。对 \(z\ne0\)，两种维数下都有
\[
 \partial_j\Phi_d(z)=-\dfrac{z_j}{\sigma_{d-1}|z|^d},
\]
它也局部可积。该式同时给出分布的一阶导数：在挖去小球后分部积分，球面边界项由 \(|\Phi_d(\varepsilon)|\varepsilon^{d-1}\) 控制，随 \(\varepsilon\downarrow0\) 趋零。

对局部 \(x\) 范围，截断核到一个固定大球即可计算卷积。局部可积核的平移在 \(L^1\) 中连续，因此核与有界 \(f\) 的卷积连续；对一阶导数核同样成立。由分布卷积的求导公式，这些连续函数是 \(Nf\) 的弱一阶导数。可用内部磨光并在坐标线段上取一致极限，说明连续弱导数就是经典导数，所以 \(Nf\in C^1\)。

最后对 \(\varphi\in C_c^\infty\) 换序，
\[
 \int Nf(x)\cdot(-\Delta\varphi(x))\,\mathrm{d}x
 =\int f(y)\left(\int\Phi_d(x-y)\cdot(-\Delta\varphi(x))\,\mathrm{d}x\right)\,\mathrm{d}y
 =\int f(y)\varphi(y)\,\mathrm{d}y.
\]
两个紧区域内的核绝对积分有限，故换序合法；内层等式是基本解的定义。
\end{proof}

势在源区之外调和。其远处行为主要由总源量 \(M=\int f\) 决定：若支集包含于 \(B(0,R_0)\)，对 \(|x|>2R_0\) 用核的中值估计，可得
\[
 Nf(x)=\dfrac{M}{(d-2)\sigma_{d-1}}|x|^{2-d}+O(|x|^{1-d})\quad(d\geq3),
\]
而二维有 \(Nf(x)=-(M/(2\pi))\log|x|+O(|x|^{-1})\)。这里常数依赖密度的一阶绝对矩及支集范围；若 \(M=0\)，首项抵消，衰减随之改善。二维非零总源量通常产生对数增长，并非无穷远趋零的势。

\ProseParagraphBreak

一个具体模型是 \(d\geq3\)、\(f=\mathbf1_{B(0,R)}\)。旋转不变性使势只依赖 \(r=|x|\)。将径向方程 \(-(r^{d-1}u')'=r^{d-1}f(r)\) 分段积分，并要求中心有限、\(u,u'\) 在 \(R\) 处连续、无穷远趋零，得到
\[
 Nf(x)=
 \begin{cases}
 \dfrac{R^2}{2(d-2)}-\dfrac{r^2}{2d},&r\leq R,\\[4pt]
 \dfrac{R^d}{d(d-2)}r^{2-d},&r\geq R.
 \end{cases}
\]
这两个表达式的一阶导数在 \(R\) 处均为 \(-R/d\)，所以没有额外的球面源；二阶导数却不连续。它展示了有界密度产生 \(C^1\) 势但未必产生 \(C^2\) 势的边界。要识别此径向构造确为 \(Nf\)，注意它与前一命题的势之差连续、分布调和且无穷远趋零，由\cref{b2:prop:13-converse-mean}和 Liouville 定理，该差只能为零。

\subsection{正则性初步}
\label{b2:subsec:130304}

为了再求一次导数，先保留 Hessian 核的角向抵消。对 \(z\ne0\)，令
\[
 K_{ij}(z)\coloneqq\dfrac1{\sigma_{d-1}}
 \left(\dfrac{d z_i z_j}{|z|^{d+2}}-\dfrac{\delta_{ij}}{|z|^d}\right).
\]
它满足 \(|K_{ij}(z)|\leq C_d|z|^{-d}\)，但在每个以原点为中心的球面上的平均为零，因为 \(\int\theta_i\theta_j\,\mathrm{d}S=(\sigma_{d-1}/d)\delta_{ij}\)。

\begin{lemma}\label{b2:lem:13-hessian-kernel}
在分布意义下，
\[
 \partial_i\partial_j\Phi_d
 =\operatorname{pv}K_{ij}-\dfrac{\delta_{ij}}d\delta_0,
\]
其中\term[zhu-zhi]{主值}[principal value]在本式中指挖去中心小球后取极限：
\[
 \langle\operatorname{pv}K_{ij},\varphi\rangle
 =\lim_{\varepsilon\downarrow0}\int_{|z|>\varepsilon}K_{ij}(z)\varphi(z)\,\mathrm{d}z.
\]
\end{lemma}
\begin{proof}
在单位球内减去 \(\varphi(0)\)，球面平均为零使被减项的截断积分为零。其余部分满足 \(|\varphi(z)-\varphi(0)|\leq C|z|\)，与核相乘后绝对可积，所以主值存在。

从 \(\partial_i\Phi_d=-z_i/(\sigma_{d-1}|z|^d)\) 出发，在小球外对 \(-\int\partial_i\Phi_d\partial_j\varphi\) 分部积分。体积项为 \(\int K_{ij}\varphi\)；内球的外法向为 \(-\theta\)，故需要减去的边界项为
\[
 \dfrac1{\sigma_{d-1}}\int_{\partial B(0,1)}\theta_i\theta_j\varphi(\varepsilon\theta)\,\mathrm{d}S_\theta
 \longrightarrow\dfrac{\delta_{ij}}d\varphi(0).
\]
这给出原点项及其负号。将 \(i=j\) 后求和，主值核的迹为零，正好恢复 \(\Delta\Phi_d=-\delta_0\)。
\end{proof}

\begin{theorem}[Hölder 密度的 Newton 势]\label{b2:thm:13-newton-holder}
若 \(f\in C_c^{0,\alpha}(\mathbb R^d)\)、\(0<\alpha\leq1\)，则 \(Nf\in C^2(\mathbb R^d)\)，且 \(-\Delta Nf=f\) 逐点成立。其二阶导数可写为
\begin{equation}\label{b2:eq:13-newton-hessian}
 \begin{aligned}
 \partial_i\partial_jNf(x)={}&
 \int_{|z|<1}K_{ij}(z)\bigl(f(x-z)-f(x)\bigr)\,\mathrm{d}z\\
 &+\int_{|z|\geq1}K_{ij}(z)f(x-z)\,\mathrm{d}z
 -\dfrac{\delta_{ij}}d f(x).
 \end{aligned}
\end{equation}
\end{theorem}
\begin{proof}
近处的绝对值由 \(C_d[f]_{C^{0,\alpha}}|z|^{\alpha-d}\) 控制，其径向积分为 \(\int_0^1r^{\alpha-1}\,\mathrm{d}r<\infty\)。远处核没有局部奇性，且对紧的 \(x\) 范围，\(f(x-z)\) 只在一个固定有界 \(z\) 区域内非零。因此两个积分都绝对收敛，并由控制收敛连续依赖 \(x\)。

前一引理与角向抵消给出该表达式作为分布二阶导数。也可先将 \(f\) 磨光，使用引理对光滑密度求导；磨光后的 Hölder 半范数受同一常数控制，上述可积界允许在局部取极限。由\cref{b2:prop:13-newton-c1}，一阶导数已连续；二阶弱导数又由连续函数给出，在坐标线段上对磨光函数的导数取一致极限，便知它们是经典二阶导数。最后对 \(i=j\) 求和，两个积分的核迹均为零，得到 \(-\Delta Nf=f\)。
\end{proof}

Hölder 条件在证明中只用于使 \(\int_0^1\omega_f(r)\,\mathrm{d}r/r\) 有限。因此同一证明也适用于具有这一 Dini 模连续度的紧支集连续密度。单纯连续性只保证 \(\omega_f(r)\to0\)，并不保证该积分有限。

\ProseParagraphBreak

确实，在二维原点附近令 \(r=|x|\)、\(\ell(r)=\log(e/r)\)，取
\[
 u(x)=(x_1^2-x_2^2)\log\ell(r),\quad u(0)=0.
\]
函数及其一阶导数在原点连续，因为它们分别为 \(O(r^2\log\ell(r))\) 和 \(O(r\log\ell(r))\)。对 \(r>0\)，极坐标计算给出
\[
 -\Delta u(x)=\left(\dfrac4{\ell(r)}+\dfrac1{\ell(r)^2}\right)\cos(2\theta),
\]
右侧在原点补零后连续。但沿 \(x_1\) 轴的二阶对称差商为 \(2\log\ell(|h|)\)，趋于无穷，所以 \(u\) 不属于 \(C^2\)。分部积分时小圆边界项趋零，故原点没有额外的点质量，等式确在分布意义下成立。

\ProseParagraphBreak

再乘一个在原点附近为 \(1\) 的光滑紧支集截断，可以将此例放到全平面，所得右端 \(f\) 仍连续且紧支集。如果 \(Nf\) 是 \(C^2\) 的，那么 \(u-Nf\) 连续且分布调和，因而光滑，就会迫使 \(u\) 为 \(C^2\)，矛盾。这个例子说明二阶奇核的抵消需要定量连续性；把“连续”直接当作“可在核下求两次导数”的许可，会遗漏实质条件。

% !TeX root = ../../Book2.tex
\section{Green 函数}
\label{b2:sec:1304}

基本解在整个空间处理点源，Green 函数还把区域和边界条件包含进来。在 Dirichlet 问题中，它由基本解减去一个调和修正，使边界值归零。修正是否存在取决于区域，不能因为形式上写出了一个差，就认为边值问题已经解决。

\subsection{有界区域}
\label{b2:subsec:130401}

\chapref{b2:ch:10}已经把 Green 函数用于一维边值算子的逆；这里明确负 Laplace 算子在有界区域上的 Dirichlet 版本。以下各项性质均以所述经典 Green 函数存在为前提，一般区域上的存在性并未因此取得。

\Needspace{6\baselineskip}
\begin{definition}\label{b2:def:13-dirichlet-green}
设整数 \(d\geq2\)，\(\Omega\subset\mathbb R^d\) 为有界连通 \(C^2\) 开集，\(\Phi_d\) 为\cref{b2:thm:13-green-representation}中明确归一化的标准径向负 Laplace 基本解。若对每个 \(x\in\Omega\)，存在实值 \(h_x\in C^2(\Omega)\cap C^1(\overline\Omega)\)，满足
\[
 \Delta_yh_x=0\quad\text{于 }\Omega,
 \quad h_x(y)=\Phi_d(x-y)\quad(y\in\partial\Omega),
\]
则定义
\[
 G_\Omega(x,y)\coloneqq\Phi_d(x-y)-h_x(y)\quad(y\ne x),
\]
称为本节经典框架下的 Dirichlet Green 函数。
\end{definition}
\symidx{\(G_\Omega(x,y)\)：Dirichlet Green 函数}

这个定义规定了点源归一化、调和修正及零边界条件。一般区域上还须证明所需的 \(h_x\) 存在；本节将在球上直接构造它，并在半空间写出对应公式。这里不预设任意粗糙边界都具有连续 Dirichlet 解。

\Needspace{6\baselineskip}
\begin{proposition}[Green 函数的唯一性与正性]\label{b2:prop:13-green-positive}
设整数 \(d\geq2\)，\(\Omega\subset\mathbb R^d\) 为有界连通 \(C^2\) 开集。假设其经典 Dirichlet Green 函数 \(G_\Omega(x,y)=\Phi_d(x-y)-h_x(y)\) 存在，其中 \(\Phi_d\) 为\cref{b2:thm:13-green-representation}中明确归一化的标准径向负 Laplace 基本解，实值 \(h_x\in C^2(\Omega)\cap C^1(\overline\Omega)\) 在 \(\Omega\) 内调和，且在边界等于 \(\Phi_d(x-\cdot)\)。则 \(G_\Omega\) 唯一，并对不同的 \(x,y\in\Omega\) 有 \(G_\Omega(x,y)>0\)。
\end{proposition}
\begin{proof}
两个候选的调和修正具有相同边界值，最大值原理使它们相等，故唯一。固定 \(x\)，当 \(y\to x\) 时，\(\Phi_d(x-y)\to+\infty\)，而 \(h_x\) 在该点附近有界。因此在足够小的内球边界上 \(G_\Omega>0\)，在外边界上为零。对挖孔区域使用最小值原理，得到 \(G_\Omega\geq0\)。当 \(d\geq2\) 时，连通开集去掉一个点仍可用绕开该点的折线连接；所以强最小值原理排除了其他内点取零的可能。
\end{proof}

\Needspace{6\baselineskip}
\begin{proposition}[Green 函数的对称性]\label{b2:prop:13-green-symmetry}
设整数 \(d\geq2\)，\(\Omega\subset\mathbb R^d\) 为有界连通 \(C^2\) 开集。假设其经典 Dirichlet Green 函数 \(G_\Omega(x,y)=\Phi_d(x-y)-h_x(y)\) 存在，其中 \(\Phi_d\) 为\cref{b2:thm:13-green-representation}中明确归一化的标准径向负 Laplace 基本解，实值 \(h_x\in C^2(\Omega)\cap C^1(\overline\Omega)\) 在 \(\Omega\) 内调和，且在边界等于 \(\Phi_d(x-\cdot)\)。则对不同的 \(x,z\in\Omega\)，有 \(G_\Omega(x,z)=G_\Omega(z,x)\)。
\end{proposition}
\begin{proof}
对函数 \(G_x(y)=G_\Omega(x,y)\)、\(G_z(y)=G_\Omega(z,y)\)，在去掉 \(x,z\) 附近两只不交小球后的区域使用第二 Green 恒等式。体积项为零，外边界项也为零，因为两个函数都在那里取零。若它们只在边界为 \(C^1\)，先在向内平移的光滑边界上使用公式再取极限：有限个 \(C^2\) 边界图表给出趋向原边界的法向与面积因子，函数及一阶导数的连续性保证边界积分收敛。

在以 \(x\) 为中心的小球上，\(G_z\) 光滑，而 \(G_x\) 的奇性与 \(\Phi_d(x-y)\) 相同。与\cref{b2:thm:13-green-representation}相同的内球计算，使这一边界的 \(G_x\partial_\nu G_z-G_z\partial_\nu G_x\) 积分趋于 \(-G_z(x)\)。在以 \(z\) 为中心的小球上，极点角色交换，极限为 \(G_x(z)\)。两项之和为零，即得对称性。
\end{proof}

因此固定 \(y\) 时，\(G_\Omega(x,y)\) 在 \(x\ne y\) 处也调和。对称性来自算子与边界条件的共同结构；若改成非自伴算子，一般应比较原算子与伴随算子的两个 Green 核，而不能直接沿用同核对称式。

\subsection{Dirichlet 问题}
\label{b2:subsec:130402}

\Needspace{6\baselineskip}
\begin{theorem}[Dirichlet Green 表示]\label{b2:thm:13-dirichlet-representation}
设整数 \(d\geq2\)，\(\Omega\subset\mathbb R^d\) 为有界连通 \(C^2\) 开集。假设其经典 Dirichlet Green 函数 \(G_\Omega(x,y)=\Phi_d(x-y)-h_x(y)\) 存在，其中 \(\Phi_d\) 为\cref{b2:thm:13-green-representation}中明确归一化的标准径向负 Laplace 基本解，实值 \(h_x\in C^2(\Omega)\cap C^1(\overline\Omega)\) 在 \(\Omega\) 内调和，且在边界等于 \(\Phi_d(x-\cdot)\)。另设 \(\mathbb K\in\{\mathbb R,\mathbb C\}\)，\(U\) 为包含 \(\overline\Omega\) 的开集，\(u\in C^2(U;\mathbb K)\)，\(f=-\Delta u\) 于 \(\Omega\)、\(g=u|_{\partial\Omega}\)。则对每个 \(x\in\Omega\)，
\begin{equation}\label{b2:eq:13-dirichlet-representation}
 u(x)=\int_\Omega G_\Omega(x,y)\cdot f(y)\,\mathrm{d}y
      +\int_{\partial\Omega}P_\Omega(x,y)\cdot g(y)\,\mathrm{d}S_y,
\end{equation}
其中 \(P_\Omega(x,y)\coloneqq-\partial_{\nu_y}G_\Omega(x,y)\)。这个边界核非负，且对每个 \(x\in\Omega\) 的边界积分为 \(1\)。
\end{theorem}
\begin{proof}
将基本解表示中的核换成 \(G_\Omega(x,\cdot)\)。两者的奇性相同，内球极限仍产生 \(u(x)\)；调和修正不增加体积源。外边界上 \(G_\Omega=0\)，原来乘以 \(\partial_\nu u\) 的项消失，留下 \(-g\partial_\nu G_\Omega\)，得到公式。对 \(C^1\) 边界迹的极限使用前一证明中的向内平行面处理。

Green 函数在内部非负、边界为零，所以向外的一阶法向导数不大于零；这是从边界向内取差商直接得到的结论。故 \(P_\Omega\geq0\)。再把 \(u=1\)、\(f=0\)、\(g=1\) 代入表示式，即得核的积分为 \(1\)。
\end{proof}

这时边界数据以加权平均的方式进入解。公式首先表示一个已有解；要把右侧直接用作存在性构造，还须检查它在内部调和、以及趋向边界时是否恢复给定的连续数据。下面的显式核将完成这些步骤。

\ProseParagraphBreak

Poisson 方程的体积项与边界项承担不同任务。前者产生内部源，后者修正边界值；若给定的是零 Dirichlet 数据，前者的 Green 核为正，所以 \(f\geq0\) 时解非负。这与一维或球上的正源例子一致，也再次核对了采用 \(-\Delta\) 后的符号。


\subsection{球的 Green 函数与 Poisson 核}
\label{b2:subsec:130403}

先看球 \(B=B(0,R)\)。对 \(0<|x|<R\)，令 \(x^*=R^2x/|x|^2\)，它位于球外。这一反演点满足：当 \(|y|=R\) 时，
\[
 |x^*-y|=\dfrac R{|x|}|x-y|.
\]
等式由两边平方后直接展开得到。它使原点源和球外像点的基本解在球面上具有可精确配平的比例。

\Needspace{6\baselineskip}
\begin{proposition}[球上的 Green 函数]\label{b2:prop:13-ball-green}
设整数 \(d\geq2\)、\(R>0\)，\(B=B_R(0)\subset\mathbb R^d\)，\(\sigma_{d-1}=|S^{d-1}|\)。固定标准径向基本解
\[
 \Phi_d(z)=
 \begin{cases}
 |z|^{2-d}/((d-2)\sigma_{d-1}),&d\geq3,\\
 -(2\pi)^{-1}\log|z|,&d=2,
 \end{cases}
 \qquad z\ne0.
\]
对 \(x,y\in B\)、\(x\ne y\)，当 \(x\ne0\) 时令 \(x^*=R^2x/|x|^2\)。当 \(d\geq3\)、\(x\ne0\) 时，球上的 Dirichlet Green 函数为
\[
 G_B(x,y)=\Phi_d(x-y)
 -\left(\dfrac R{|x|}\right)^{d-2}\Phi_d(x^*-y).
\]
当 \(d=2\)、\(x\ne0\) 时，
\[
 G_B(x,y)=\dfrac1{2\pi}
 \log\left(\dfrac{|x|\cdot|x^*-y|}{R|x-y|}\right).
\]
中心处统一取 \(G_B(0,y)=\Phi_d(y)-\Phi_d(Re_1)\)。
\end{proposition}
\begin{proof}
像点 \(x^*\) 在球外，因此相应的基本解在闭球附近光滑且调和。高维情况下，将球面距离恒等式代入 \(|z|^{2-d}\)，两个边界项恰好抵消；二维情况下，同一恒等式使对数的真数为 \(1\)，边界值也为零。原点源的系数仍为 \(1\)，故原点奇性未改，满足定义中的全部要求。对 \(x=0\)，修正直接取常数，使边界归零即可。唯一性已由前文证明。
\end{proof}

\Needspace{6\baselineskip}
\begin{proposition}[球上的 Poisson 核]\label{b2:prop:13-ball-poisson}
设整数 \(d\geq2\)、\(R>0\)，\(B=B_R(0)\subset\mathbb R^d\)，\(\sigma_{d-1}=|S^{d-1}|\)。球的 Dirichlet Poisson 核为
\begin{equation}\label{b2:eq:13-ball-poisson}
 P_B(x,y)=\dfrac{R^2-|x|^2}{\sigma_{d-1}R|x-y|^d}
 \quad(|x|<R,\ |y|=R).
\end{equation}
它关于 \(x\) 调和、严格为正，且 \(\int_{\partial B}P_B(x,y)\,\mathrm{d}S_y=1\)。
\end{proposition}
\begin{proof}
在球面上 \(\nu_y=y/R\)。基本解的一阶导数给出
\[
 \partial_{\nu_y}\Phi_d(x-y)
 =-\dfrac{R^2-x\cdot y}{\sigma_{d-1}R|x-y|^d}.
\]
对像点项同样求导，并用 \(|x^*-y|=(R/|x|)|x-y|\)。其系数与分母的合并因子为 \(|x|^2/R^2\)，于是相减后的分子为
\[
 (R^2-x\cdot y)-\dfrac{|x|^2}{R^2}(R^2-x^*\cdot y)=R^2-|x|^2.
\]
这给出 \(-\partial_{\nu_y}G_B\) 的公式。二维对数中的附加常数不影响 \(y\) 导数；中心情形由直接计算或连续极限得到。

正性由公式可见，积分为 \(1\) 可把常函数代入已证明的 Green 表示式。至于 \(x\) 的调和性，对固定 \(|y|=R\) 直接使用乘积法则，得到
\[
 \Delta_x\bigl((R^2-|x|^2)|x-y|^{-d}\bigr)
 =2d(R^2-|y|^2)|x-y|^{-d-2}=0.
\]
所需导数在每个内部紧集上都没有奇性。
\end{proof}
\symidx{\(P_B(x,y)\)：球的 Poisson 核}

\subsection{球上的 Dirichlet 与 Poisson 问题}
\label{b2:subsec:13-ball-boundary}

已经构造的核还须接受两项检验：它应在内部给出调和函数，并在趋向边界时取得给定数据。前面的表示式针对已有解，下面才用显式球核完成连续边值问题的存在性。

\Needspace{6\baselineskip}
\begin{theorem}[球上的连续 Dirichlet 问题]\label{b2:thm:13-ball-dirichlet}
设整数 \(d\geq2\)、\(R>0\)，\(B=B_R(0)\subset\mathbb R^d\)，\(\sigma_{d-1}=|S^{d-1}|\)，\(\mathbb K\in\{\mathbb R,\mathbb C\}\)。令 \(P_B(x,y)=(R^2-|x|^2)/(\sigma_{d-1}R|x-y|^d)\)。对每个 \(g\in C(\partial B;\mathbb K)\)，函数
\[
 u(x)=\int_{\partial B}P_B(x,y)\cdot g(y)\,\mathrm{d}S_y\quad(x\in B)
\]
与边界上的 \(u=g\) 共同定义了 \(C(\overline B)\) 函数。它在 \(B\) 内光滑且调和，并且是这一函数类中唯一满足该边界数据的解。
\end{theorem}
\begin{proof}
内部紧集离边界有正距离，核的任意阶 \(x\) 导数在该紧集乘球面上统一有界，所以可在积分下逐项求导，得到内部光滑性与调和性。

固定边界点 \(\zeta\)。给定 \(\varepsilon>0\)，由 \(g\) 连续性选 \(\delta>0\)，使 \(|g(y)-g(\zeta)|<\varepsilon\) 当 \(|y-\zeta|<\delta\)。利用正性与总质量为 \(1\)，
\[
 |u(x)-g(\zeta)|
 \leq\varepsilon+2\|g\|_\infty
       \int_{|y-\zeta|\geq\delta}P_B(x,y)\,\mathrm{d}S_y.
\]
当 \(x\to\zeta\) 且 \(|x-\zeta|<\delta/2\) 时，远处有 \(|x-y|\geq\delta/2\)，而分子 \(R^2-|x|^2\to0\)。所以尾积分趋零，得到边界连续性。唯一性由有界域上的最大值原理取得。
\end{proof}

现在用球上的连续边值解证明前面延后的局部均值逆命题。

\begin{proof}[\cref{b2:prop:13-local-converse-mean}的证明]
先处理实值函数。取任意 \(\overline B=\overline{B(a,R)}\subset\Omega\)。当 \(d\geq2\) 时，用\cref{b2:thm:13-ball-dirichlet}（将球心平移到 \(a\)）取得 \(v\in C(\overline B)\)，它在 \(B\) 内调和，边界值为 \(u\)；当 \(d=1\) 时，直接以两端值的仿射插值作为 \(v\)。差 \(w=u-v\) 连续，在边界为零，并在每个内部点满足所有充分小半径的球面平均性质。

若 \(M=\max_{\overline B}w>0\)，最大值只能在内部达到。集合 \(F=\{x\in B:w(x)=M\}\) 非空且在 \(B\) 中闭。对任意 \(x\in F\)，取同时满足局部均值假设且闭球仍含于 \(B\) 的小半径。函数 \(M-w\geq0\) 在每个这样的小球面上的积分为零；连续性使它在整个球面上为零。让半径遍历一个小区间，就得到 \(w=M\) 在 \(x\) 的整个邻域成立。因此 \(F\) 在 \(B\) 中也开。球的连通性推出 \(F=B\)，再趋向边界便与 \(w=0\) 矛盾。所以 \(w\leq0\)；对 \(-w\) 重复论证得到 \(w\geq0\)，从而 \(u=v\) 于 \(B\)。

所取内部闭球任意，而 \(u\) 在各个内部球上都等于光滑调和函数 \(v\)，故 \(u\) 在 \(\Omega\) 内光滑且调和。复值情形分别对实部、虚部使用相同论证。
\end{proof}

例如单位圆上的边界数据 \(g(e^{\mathrm{i}\theta})=\cos(m\theta)\) 对应于圆盘内 \(u(r,\theta)=r^m\cos(m\theta)\)。可直接在极坐标下检查调和性和边界值，再由唯一性确认它就是 Poisson 积分。高频边界振荡在内部被 \(r^m\) 压低；当 \(r<1\) 固定时，频率越高，影响越小。这也与\chapref{b2:ch:02}圆周 Poisson 平均的系数权重一致。

\Needspace{6\baselineskip}
\begin{corollary}[球上的 Poisson 方程]\label{b2:cor:13-ball-poisson-problem}
设整数 \(d\geq2\)、\(R>0\)，\(B=B_R(0)\subset\mathbb R^d\)，\(\mathbb K\in\{\mathbb R,\mathbb C\}\)。设 \(0<\alpha\leq1\)、\(f\in C^{0,\alpha}(\overline B;\mathbb K)\)，\(g\in C(\partial B;\mathbb K)\)。则存在唯一的 \(u\in C^2(B;\mathbb K)\cap C(\overline B;\mathbb K)\)，满足 \(-\Delta u=f\)、\(u|_{\partial B}=g\)。
\end{corollary}
\begin{proof}
先将 \(f\) 延拓为某个 \(F\in C_c^{0,\alpha}(\mathbb R^d)\)。实值情形可用
\[
 F_0(x)=\inf_{y\in\overline B}\bigl(f(y)+L|x-y|^\alpha\bigr),
\]
其中 \(L\) 为 Hölder 常数。由 \(|x-z|^\alpha\leq|x-y|^\alpha+|y-z|^\alpha\)，可证 \(F_0\) 保持该 Hölder 常数，且在 \(\overline B\) 上等于 \(f\)。乘以在闭球附近为 \(1\) 的光滑紧支集截断，得到 \(F\)；复值时按实、虚部分别延拓。

令 \(v=NF\)，由\cref{b2:thm:13-newton-holder}知 \(v\in C^2\)、\(-\Delta v=f\) 于球内。再对连续边界值 \(g-v|_{\partial B}\) 使用 Poisson 积分，所得调和函数记为 \(h\)。于是 \(u=v+h\) 满足方程与边界数据。两个候选之差调和且在边界为零，故唯一；这也说明构造不依赖于所选的延拓。
\end{proof}


\subsection{半空间的 Green 函数与 Dirichlet 问题}
\label{b2:subsec:13-halfspace-boundary}

再看上半空间 \(H=\{(x',t)\in\mathbb R^{d-1}\times\mathbb R:t>0\}\)。固定 \(x=(x',t)\)，其关于边界平面的反射点为 \(x^*=(x',-t)\)。令
\[
 G_H(x,y)=\Phi_d(x-y)-\Phi_d(x^*-y).
\]
反射点在区域外，且平面上的点到 \(x,x^*\) 等距，所以核在边界为零。对 \(y\) 而言，区域内只有原极点，满足 \(-\Delta_yG_H(x,\cdot)=\delta_x\)。这里没有把有界区域上的唯一性直接移到半空间；无穷远的条件还需单独指定。

\Needspace{6\baselineskip}
\begin{proposition}[半空间的 Poisson 核]\label{b2:prop:13-halfspace-poisson}
设整数 \(d\geq2\)，\(H=\{(x',t)\in\mathbb R^{d-1}\times\mathbb R:t>0\}\)，\(\sigma_{d-1}=|S^{d-1}|\)。对 \(x=(x',t)\in H\)、\(y'\in\mathbb R^{d-1}\)，半空间负 Laplace 算子的 Dirichlet Poisson 核为
\[
 P_H(x,y')=\dfrac2{\sigma_{d-1}}
       \dfrac{t}{\bigl(|x'-y'|^2+t^2\bigr)^{d/2}}.
\]
它严格为正，对 \(x\) 调和，并有 \(\int_{\mathbb R^{d-1}}P_H(x,y')\,\mathrm{d}y'=1\)。
\end{proposition}
\begin{proof}
边界外法向为 \(-e_d\)，因此 \(-\partial_{\nu_y}G_H=\partial_{y_d}G_H\)。在 \(y_d=0\) 处，两个极点项的竖直导数大小相同、符号相反，相减后得到系数 \(2t/\sigma_{d-1}\)。对 \(x\) 的调和性可直接对有理核求导验证；也可注意它是原点外调和基本解的一个坐标导数。

尺度换元使总积分不依赖 \(t\)。为计算常数，令
\[
 I=\int_{\mathbb R^{d-1}}(1+|z|^2)^{-d/2}\,\mathrm{d}z,
 \quad J=\int_0^\infty e^{-\pi s^2}s^{d-1}\,\mathrm{d}s.
\]
在 \(d\) 维 Gauss 积分的上半空间内作换元 \((x',x_d)=(sz,s)\)，Jacobi 因子为 \(s^{d-1}\)，再缩放径向变量，得到 \(1/2=IJ\)。而整个空间的极坐标公式给出 \(1=\sigma_{d-1}J\)，因此 \(I=\sigma_{d-1}/2\)，所写核的积分为 \(1\)。
\end{proof}
\symidx{\(P_H(x,y')\)：半空间 Poisson 核}

这正是\secref{b2:sec:0304}中定义在 \(\mathbb R^{d-1}\) 上的 Poisson 逼近恒等。那里先用积分定义归一化常数，这里通过边界几何得到了它的具体数值。

\Needspace{6\baselineskip}
\begin{theorem}[半空间的有界连续 Dirichlet 问题]\label{b2:thm:13-halfspace-dirichlet}
设整数 \(d\geq2\)，\(H=\{(x',t)\in\mathbb R^{d-1}\times\mathbb R:t>0\}\)，\(\sigma_{d-1}=|S^{d-1}|\)，\(\mathbb K\in\{\mathbb R,\mathbb C\}\)。令 \(P_H((x',t),y')=2t/[\sigma_{d-1}(|x'-y'|^2+t^2)^{d/2}]\)。若 \(g\in C_b(\mathbb R^{d-1};\mathbb K)\)，则
\[
 u(x',t)=\int_{\mathbb R^{d-1}}P_H((x',t),y')\cdot g(y')\,\mathrm{d}y'
\]
是在 \(H\) 内调和、在 \(\overline H\) 上连续且有界、边界值为 \(g\) 的唯一函数。并且 \(\|u\|_\infty\leq\|g\|_\infty\)。
\end{theorem}
\begin{proof}
核非负且积分为 \(1\)，所以积分绝对收敛并给出范数界。对每个紧正高度区间，各阶内部导数核在边界变量中具有统一的可积尾部，因此可在积分下求导，得到内部调和性与光滑性。

固定边界点 \((a',0)\)，按 \(|y'-a'|<\delta\) 与其余部分分割误差。近处由 \(g\) 在该点连续而小；当 \(|x'-a'|<\delta/2\) 时，远处的核积分不超过
\[
 C_d t\int_{\delta/2}^\infty r^{-d}r^{d-2}\,\mathrm{d}r
 \leq C_d\,t/\delta,
\]
所以 \(x'\to a'\)、\(t\downarrow0\) 时误差趋零，得到边界连续性。

唯一性须使用有界条件。若 \(w\) 为两个解之差，则其边界值为零，在平面另一侧作奇反射，得到整个 \(\mathbb R^d\) 上的有界连续函数 \(\widetilde w\)。它在平面两侧分别调和；在平面上的每个点，以该点为中心的球面平均因反射抵消而为零，等于函数值。因此每个点都有足够小半径的球面平均性质。

由已经完成证明的\cref{b2:prop:13-local-converse-mean}，\(\widetilde w\) 在整个空间光滑且调和。它又全局有界，\cref{b2:thm:13-liouville}使其为常数；奇对称性及平面上的零值迫使该常数为零。因此两个半空间解相同。
\end{proof}

如果去掉有界性，函数 \(u(x',t)=t\) 就是非零的零边界调和函数，唯一性失败。对一般无界区域，边界值只控制有限位置的边界，不能替代无穷远增长条件。另一方面，若 \(g\) 有界且一致连续，则由逼近恒等理论还可得到 \(u(\cdot,t)\to g\) 的一致收敛；只有逐点连续时，本定理只断言相应的边界连续性。

\subsection{镜像法}
\label{b2:subsec:130404}

在半空间中，补入符号相反的反射点源，使平面上两项恰好抵消；在球中，则用反演点、维数相关的系数和二维的对数常数完成同样的边界配平。这类构造称为\term[jingxiang-fa]{镜像法}[method of images]。两种像点的几何关系见\cref{b2:fig:13-images}。
% 本书原创：平面反射与球面反演的二维截面示意。
\begin{figure}[htbp]
\centering
\begin{tikzpicture}[x=1cm,y=1cm,>=stealth,every node/.style={font=\small},line width=.65pt]
\begin{scope}
\fill[AnalysisBlue!7] (0,0) circle(1.55);
\draw[AnalysisBlue] (0,0) circle(1.55);
\draw[->,gray] (-1.85,0)--(3.6,0);
\fill (0,0) circle(1.2pt) node[below left]{\(0\)};
\fill[AnalysisBlue] (.6975,0) circle(2pt) node[below]{\(x\)};
\fill[AnalysisBlue] (3.4444,0) circle(2pt) node[below]{\(x^*\)};
\fill (0.775,1.34234) circle(1.5pt) node[above]{\(y\)};
\draw[dashed,gray] (.6975,0)--(.775,1.34234)--(3.4444,0);
\node at (.6,2.1){球面反演};
\node at (.5,-2.0){\(R=1,\quad |x|=0.45,\quad |x^*|=1/0.45\)};
\end{scope}
\begin{scope}[xshift=8cm]
\fill[AnalysisBlue!7] (-1.7,0) rectangle(2.2,1.65);
\draw[gray] (-1.85,0)--(2.35,0) node[right]{边界};
\draw[dashed,gray] (0,-1.2)--(0,1.2);
\fill[AnalysisBlue] (0,1.2) circle(2pt) node[above]{\(x=(x',t)\)};
\fill[AnalysisBlue] (0,-1.2) circle(2pt) node[below]{\(x^*=(x',-t)\)};
\fill (1.8,0) circle(1.5pt) node[above right]{\(y\)};
\draw[dashed,gray] (0,1.2)--(1.8,0)--(0,-1.2);
\node at (.25,2.1){平面反射};
\node at (.25,-2.0){\(|x-y|=|x^*-y|\)};
\end{scope}
\end{tikzpicture}
\caption[边界配平中的两种像点]{边界配平中的两种像点。左图是单位球的二维截面，球面距离相差固定比例；右图的边界是平面，两条距离相等。像点均位于区域外，基本解的系数须分别按球面反演与平面反射公式选取。}
\label{b2:fig:13-images}
\end{figure}


镜像的几何位置与系数各有职责。平面反射保持到边界点的距离，因此系数为 \(1\) 即可；球面反演只保持一个固定距离比，因而高维幂核必须乘上相应幂次，二维还要减去常数。若遗漏这些系数，所得表达式可能在内部仍满足点源方程，却不满足零边界条件。

\ProseParagraphBreak

两个互相垂直的平面可以通过有限次反射配出更多像点；章末习题给出第一象限的四点构造，并要求检查每条边界及正性。一般曲面则没有这样的有限反射关系。此时不能把任意一个“对称点”放到区域外就称为 Green 函数，而应回到调和修正的边界问题，或使用变分和势论方法。

\ProseParagraphBreak

本章给出的球与半空间公式分别提供了有界区域与无界区域的基本模型。它们都从同一个基本解出发，但边界配平、核的归一化以及唯一性所需的函数类各不相同。进入一般椭圆边值问题时，正是这些需要逐项保留的条件决定后面的理论能够推广到哪里。


\begin{chaptersummary}
均值性质把调和函数的局部微分方程变成平均约束，由此得到内部极值的排除、正函数的比例比较及任意阶局部估计。Green 恒等式同时记录内部源与边界通量，区分了三类边界条件的相容性和唯一性。无界区域还必须限制增长，不能只凭有限位置的边界值断言解唯一。

\ProseParagraphBreak

Newton 势先给出全空间特解；一阶导数的核局部可积，二阶导数则要保留角向抵消和原点项。Hölder 或 Dini 条件提供了这种抵消后的可积控制，单纯连续右端不普遍保证 \(C^2\) 正则性。Green 函数再通过调和修正满足边界，球和半空间的显式构造展示了表示、边界极限与唯一性如何闭合。下一编将用函数空间和变分方法处理无法直接写出核的区域与算子。
\end{chaptersummary}

\BookExercises

前几题组合平均、核与能量方法；后面的题目改变区域、边界连续性或无穷远条件，用以检验本章各项假设的作用。

% 自拟：二次型边界数据的调和修正，不只辨认一个调和多项式。
\begin{exercise}\label{b2:ex:13-quadratic-boundary}
设 \(A\) 为实对称 \(d\times d\) 矩阵。在 \(B(0,R)\) 上求边界数据 \(g(y)=y^TAy\) 的调和延拓，并求其中心值。若 \(A\) 半正定，说明延拓在球内为何非负。
\end{exercise}
\begin{solution}
因为 \(\Delta(x^TAx)=2\operatorname{tr}A\)，而 \(\Delta|x|^2=2d\)，函数
\[
 u(x)=x^TAx+\dfrac{\operatorname{tr}A}{d}(R^2-|x|^2)
\]
调和，在边界上等于 \(g\)。由 Dirichlet 唯一性，它就是所求延拓，且 \(u(0)=R^2\operatorname{tr}A/d\)。半正定时边界数据非负，最小值原理保证内部非负；这项判断不需要从修正公式中逐项比较符号。
\end{solution}

% 改编：Axler2001Harmonic, Exercises3,2（印刷56页）的取等问题；
% 增加一般半径的推导与不连续边界极值对象的辨析。
\begin{exercise}\label{b2:ex:13-sharp-harnack}
设 \(d\geq2\)，\(u>0\) 在 \(B(0,R)\) 调和，\(|x|=r<R\)、\(s=r/R\)。证明
\[
 \dfrac{1-s}{(1+s)^{d-1}}u(0)
 \leq u(x)\leq
 \dfrac{1+s}{(1-s)^{d-1}}u(0).
\]
当 \(x\ne0\) 时，分别给出取等的正调和函数；说明这些例子为何没有连续延拓到整个球面。
\end{exercise}
\begin{solution}
先假设函数在闭球附近调和。由 Poisson 表示、\(R-r\leq|x-y|\leq R+r\) 以及球面平均，
\[
 \dfrac{R^{d-2}(R^2-r^2)}{(R+r)^d}u(0)
 \leq u(x)\leq
 \dfrac{R^{d-2}(R^2-r^2)}{(R-r)^d}u(0),
\]
化简即得。一般情形在半径 \(\rho\in(r,R)\) 的球上应用，再令 \(\rho\uparrow R\)。

固定边界点 \(\zeta\)，函数 \(u_\zeta(z)=\sigma_{d-1}R^{d-1}P_B(z,\zeta)\) 在开球内为正且调和，中心值为 \(1\)。当 \(\zeta=Rx/r\) 时，在 \(x\) 处取得上界；当 \(\zeta=-Rx/r\) 时取得下界。沿着趋向 \(\zeta\) 的径向路径，Poisson 核具有边界奇性，所以这些函数没有连续延拓到整个闭球；它们仍符合题目的开球假设。
\end{solution}

% 自拟：零总源量后的偶极项，使用已建立Hessian核作远场余项估计。
\begin{exercise}\label{b2:ex:13-dipole}
设 \(d\geq2\)，\(f\) 有界且紧支集，\(\int f=0\)，记 \(m=\int y f(y)\,\mathrm{d}y\)。证明
\[
 Nf(x)=\dfrac{x\cdot m}{\sigma_{d-1}|x|^d}+O(|x|^{-d})
 \quad(|x|\to\infty).
\]
构造 \(m\ne0\) 的光滑紧支集密度，并说明当 \(m=0\) 时衰减如何改善。
\end{exercise}
\begin{solution}
设支集包含于 \(B(0,R_0)\)。当 \(|x|>2R_0\) 时，对 \(\Phi_d(x-y)\) 在 \(y=0\) 处作一次 Taylor 展开。线段上二阶核的模不超过 \(C_d|x|^{-d}\)，所以余项积分由 \(C_d|x|^{-d}\int|y|^2|f(y)|\,\mathrm{d}y\) 控制。常数项因零总质量消失，线性项为 \(-\nabla\Phi_d(x)\cdot m=x\cdot m/(\sigma_{d-1}|x|^d)\)。

取偶的非负单位积分磨光函数 \(\rho\)，以及 \(a\ne0\)，令 \(f(y)=\rho(y-a)-\rho(y+a)\)。换元给出总质量为零，第一矩为 \(2a\)，故偶极项非零。如果第一矩也为零，则只剩 \(O(|x|^{-d})\) 的余项；这是一项由矩相消产生的改进。
\end{solution}

% 自拟：有两个边界分支时的外法向与Neumann相容性。
\begin{exercise}\label{b2:ex:13-annulus-neumann}
在环域 \(a<|x|<b\)、\(d\geq2\)，要求调和函数在内、外球面的外法向导数分别为常数 \(g_a,g_b\)。给出存在解的充要条件，写出所有径向解，并说明这些解是否已覆盖全部经典解。
\end{exercise}
\begin{solution}
相容条件为 \(a^{d-1}g_a+b^{d-1}g_b=0\)。径向调和方程给出 \(u'(r)=Cr^{1-d}\)，而内边界的外法向为 \(-e_r\)，故
\[
 C=-a^{d-1}g_a=b^{d-1}g_b.
\]
条件成立时，可取 \(u(r)=C\log r+B\)（\(d=2\)），或 \(u(r)=C r^{2-d}/(2-d)+B\)（\(d\geq3\)），所以条件也充分。任意经典解与此显式解具有相同 Neumann 数据，其差由能量唯一性为常数，因此这些径向解已经给出全部解。再指定环域上的平均值即可固定 \(B\)。
\end{solution}

% 自拟：同时改变源项与边界值，给出可取等的稳定性常数。
\begin{exercise}\label{b2:ex:13-source-stability}
设 \(u,v\in C^2(B(0,R))\cap C(\overline{B(0,R)})\)，且 \(-\Delta u=f\)、\(-\Delta v=k\)，其中 \(f-k\) 有界。证明
\[
 \|u-v\|_\infty\leq
 \|u-v\|_{C(\partial B)}+\dfrac{R^2}{2d}\|f-k\|_\infty,
\]
并证明两个常数可同时取等。
\end{exercise}
\begin{solution}
记 \(M=\|u-v\|_{C(\partial B)}\)、\(F=\|f-k\|_\infty\)，令 \(w=(R^2-|x|^2)/(2d)\)，则 \(-\Delta w=1\)、边界值为零。函数 \(u-v-M-Fw\) 的 Laplace 非负，边界值不大于零，由\cref{b2:prop:13-subharmonic-comparison}得 \(u-v\leq M+Fw\)。对负差同样处理，得到绝对值界，再使用 \(\max w=R^2/(2d)\)。取 \(v=0\)、\(u=M+Fw\)、\(M,F\geq0\)，在球心即取等。
\end{solution}

% 自拟：改变区域而保持点源，比较Green核并推出正密度势的单调性。
\begin{exercise}\label{b2:ex:13-domain-monotonicity}
设 \(\overline{\Omega_1}\subset\Omega_2\)，两者满足本章经典 Green 函数的假设。证明对 \(x,y\in\Omega_1\)、\(x\ne y\)，有 \(G_{\Omega_1}(x,y)\leq G_{\Omega_2}(x,y)\)。若非负 \(f\) 紧支撑于 \(\Omega_1\)，比较两个 Green 核与 \(f\) 的积分。
\end{exercise}
\begin{solution}
固定 \(x\)。两个核之差的极点奇性相消，故可在整个 \(\Omega_1\) 上延拓为调和函数。在 \(\partial\Omega_1\) 上，差等于 \(G_{\Omega_2}(x,\cdot)\geq0\)，所以最小值原理给出差非负。乘以非负密度 \(f\) 再积分便保持不等式；核在极点局部可积，若 \(f\) 有界或相应积分有限，比较即在通常积分意义下成立。一般非负密度可先截断，再由单调收敛保留允许无穷值的比较。
\end{solution}

% 自拟：边界跳跃处依赖趋近路径，检验连续边值假设。
\begin{exercise}\label{b2:ex:13-halfplane-step}
在上半平面取边界数据 \(g(y)=\mathbf1_{(0,\infty)}(y)\)。计算其 Poisson 积分，求它在每个非零边界点的极限，以及沿 \(x=ct\)、\(t\downarrow0\) 趋于原点的极限。说明为何不存在在闭半平面连续且取该边界数据的函数。
\end{exercise}
\begin{solution}
直接积分得到
\[
 u(x,t)=\dfrac1\pi\int_0^\infty\dfrac{t}{(x-y)^2+t^2}\,\mathrm{d}y
 =\dfrac12+\dfrac1\pi\arctan(x/t).
\]
在正边界点极限为 \(1\)，负边界点为 \(0\)。沿 \(x=ct\) 的极限为 \(1/2+\pi^{-1}\arctan c\)，随 \(c\) 改变，所以原点处没有与路径无关的连续极限。题目的边界数据本身有跳跃，也已排除闭半平面上的连续延拓。内部 Poisson 积分仍是有界调和函数。
\end{solution}

% 自拟：两次反射形成四个像点，并用代数证明正性。
\begin{optionalexercise}\label{b2:ex:13-quadrant-images}
令 \(Q=\{(x_1,x_2):x_1>0,x_2>0\}\)，\(R_1,R_2\) 分别反射第一、第二坐标。对 \(x,y\in Q\)、\(x\ne y\)，验证
\[
 G(x,y)=\Phi_2(x-y)-\Phi_2(x-R_1y)-\Phi_2(x-R_2y)+\Phi_2(x-R_1R_2y)
\]
关于 \(x\) 满足单位点源方程，在两条坐标轴上为零，并严格为正。为何仅有这些方程和边界条件仍不能唯一确定无界区域上的核？
\end{optionalexercise}
\begin{solution}
三个像点都在区域外，故区域内只有 \(y\) 贡献 \(-\Delta_xG=\delta_y\)。在任一坐标轴上，等距的两对项分别抵消。若 \(x=(a,b)\)、\(y=(c,e)\)，令
\[
 A=(a-c)^2+(b-e)^2,\quad p=4ac>0,\quad q=4be>0,
\]
则四个距离平方分别为 \(A,A+p,A+q,A+p+q\)。所以
\[
 G(x,y)=\dfrac1{4\pi}\log
 \left(\dfrac{(A+p)(A+q)}{A(A+p+q)}\right)>0,
\]
因为分子比分母大 \(pq\)。但对固定 \(y\)，还可给 \(G\) 加上任意常数倍的调和函数 \(x_1x_2\)，不改变方程及坐标轴边界值。需要无穷远增长条件才能排除这种自由度；显式四像点核的远处衰减正是一种选择。
\end{solution}
